This index records the disposition of the requested manuscript rows and numbered targets. ``Proved under hypotheses'' refers only to the setting specified in the cited statement; an accompanying ``Open'' retains the broader literal path assertion. In particular, the comparison problem B.8 is not resolved by a theorem about the restricted hammock category or by a rectification theorem for diagram categories.

\begingroup
\small
\setlength{\tabcolsep}{3mm}
\renewcommand{\arraystretch}{1.12}
\newcommand{\auditref}[1]{\hyperref[res:#1]{#1}}
% The three text columns total 144 mm; four internal tabcolseps add 12 mm.
\begin{longtable}{@{}>{\raggedright\arraybackslash}p{25mm}>{\raggedright\arraybackslash}p{25mm}>{\raggedright\arraybackslash}p{94mm}@{}}
\toprule
Target & Disposition & Exact scope and numbered location \\
\midrule
\endfirsthead
\toprule
Target & Disposition & Exact scope and numbered location \\
\midrule
\endhead
\midrule
\multicolumn{3}{r@{}}{\footnotesize Continued on the following page.}\\
\endfoot
\bottomrule
\endlastfoot
\multicolumn{3}{@{}l@{}}{\textit{Chapter IX: composition, pointed structure, and linear structure}}\\
\addlinespace[2pt]
IX-037 & Proved &
Strict all-type composition and its coherent enrichment are \auditref{C.1}; the specified sequential operation is associative and unital in $\Hot$ by \auditref{C.P4}. Its localization interpretation is proved for weak verticals by \auditref{B.2}, and for the all-arrows loop model by \auditref{E.2}. General $J$ remains \auditref{B.8}.\\
\addlinespace[3pt]
IX-040 & Known under comparison; Open otherwise &
\auditref{C.5} identifies Quillen's groupoid with the fundamental groupoid of a derived function complex, for cofibrant source and fibrant target or specified replacements. Identification with the literal full path model requires the endpoint-compatible comparison retained in \auditref{B.8}. The conclusion is an equivalence of groupoids, with specified vertices.\\
\addlinespace[3pt]
IX-041 & Known under comparison; Open otherwise &
Derived mapping groupoids include all zigzag components, as explained in \auditref{C.5}; \auditref{B.2} supplies the restricted hammock realization. The corresponding identification of every literal arbitrary-vertical path vertex, with its naturality, retains the comparison requirement. \auditref{A.P1} establishes the component statement unconditionally.\\
\addlinespace[3pt]
IX-044 & Proved under hypotheses; Open generally &
\auditref{C.3} gives the smash-product correction and unit $S^0$. \auditref{C.4} constructs strictly absorbing pointed composition on cofibrant-fibrant objects of a pointed simplicial model category. The specified stabilized-path lift, with its absorption coherences, remains \auditref{C.10}.\\
\addlinespace[3pt]
IX-049 & Proved under hypotheses; Open generally &
Associative unital smash composition is proved in the pointed model of \auditref{C.4}. The exact categorical path pairings, including the proposed stronger assertion in pointed $\Cat\,\mathrm{hot}$, remain \auditref{C.10}; an unpointed equivalence supplies neither the absorption homotopies nor their compatibility.\\
\addlinespace[3pt]
IX-052 & Open for the literal ind-system &
\auditref{C.10} retains the construction in $\operatorname{Ind}(\Hot_*)$ and $\operatorname{Ind}(\Cat\,\mathrm{hot}_*)$ with all pointed transition coherences. The sufficient pointed enrichment of \auditref{C.4} is not asserted to identify these particular ind-objects.\\
\addlinespace[3pt]
IX-061 & Proved for the dg model; Open generally &
\auditref{C.8} constructs $\tau_{\leq0}\RHom_k(a,b)\in D^{\leq0}(k)$ for complexes and the stated dg enrichments. The unaxiomatized additive relative-category proposal, and its comparison with literal full path nerves, remain \auditref{C.9}.\\
\addlinespace[3pt]
IX-065 & Proved for the dg model; Open generally &
The derived-tensor composition and units, including truncation and the Koszul factor-order convention, are proved in \auditref{C.8}. Ordinary $k$-linearity of the localization alone is not the enhancement hypothesis; the broader claim is retained in \auditref{C.9}.\\
\addlinespace[3pt]
IX-070 & Proved for the dg model; Open generally &
\auditref{C.8} supplies the strong bilinear composition in its specified derived setting. \auditref{C.9} records the missing general enhancement and literal-path comparison, including the scope of the stronger additive model-category sufficient conditions.\\
\addlinespace[3pt]
IX-072 & Known under pointed comparison; Open otherwise &
\auditref{C.5} identifies the group at the zero map with Quillen's group after the stated replacements. Applying this to the manuscript's pointed full path object additionally needs a comparison carrying that basepoint, or a specified basepoint path; \auditref{C.10} retains that transfer problem.\\
\addlinespace[3pt]
\multicolumn{3}{@{}l@{}}{\textit{Chapter IX: the tests and loop--suspension assertions}}\\
\addlinespace[2pt]
IX-080 & Refuted in the weak reading; Open in the stronger reading &
\auditref{D.1} distinguishes both formulations of (b). Weak equivalence to a Thomason-fibrant category is automatic, whereas restriction-invariance can fail by \auditref{D.5}. The suitably bounded $\Cat\,\mathrm{hot}$ local-representability interpretation is retained in \auditref{D.22}, item 5. The unrestricted small-index restriction is \auditref{D.K4}.\\
\addlinespace[3pt]
IX-081 & Refuted as a completeness criterion; precise criteria qualified &
The finite example satisfies the tests and is rectifiable: \auditref{D.2}. Its missing finite products are explained by \auditref{D.3}, so those tests do not imply model or derivator completeness. \auditref{D.15} gives the family-level Kaledin criterion on its actual shapes; the smaller-shape recognition questions remain \auditref{D.22}.\\
\addlinespace[3pt]
IX-084 & Proved under hypotheses; weaker inference Refuted &
\auditref{C.6} and \auditref{D.6} prove the pointed mapping-space identities when loops and suspension are the requisite homotopy pullback and pushout. \auditref{C.7} refutes the inference from representing only $\pi_1$. Transfer to the exact pointed path model remains \auditref{C.10}.\\
\addlinespace[3pt]
IX-087 & Proved for pointed derivators; conditional for a naive family &
\auditref{C.6} proves the derivator adjunction and preservation of homotopy limits by loops and homotopy colimits by suspension. \auditref{D.7}, \auditref{D.17}, and \auditref{D.18} give sufficient diagram theories with their size hypotheses. No such structure is inferred merely from the notation $M^A[W_A^{-1}]$.\\
\addlinespace[3pt]
IX-088 & Open for the literal implication; Proved in sufficient settings &
The implication from the manuscript's path-object (c) to its restriction-invariance (b) remains \auditref{D.22}, item 5. \auditref{D.4} proves rectification implies (b), and \auditref{D.6} proves (c) implies the corresponding part of (a). These distinct implications do not establish the missing general one.\\
\addlinespace[3pt]
\multicolumn{3}{@{}l@{}}{\textit{Chapters V and VII: sequential models and endpoint properties}}\\
\addlinespace[2pt]
VIIc-A046 & Refuted in $\Cat\,\mathrm{hot}$; Proved in $\Hot$ for the path systems &
\auditref{C.P5} gives a four-object crown for which the actual $q_L$ has no inverse modulo finite zigzags of natural transformations. The specified cofinal type-system comparisons are weak equivalences by \auditref{C.P3}. \auditref{A.9} separately checks the finite subdivision warning.\\
\addlinespace[3pt]
V-I066 & Refuted in the stronger category; qualified positive comparison &
For the specified cofinal projective type systems, canonical transitive comparison of unions in $\Hot$ is \auditref{C.P3}. Promotion to $\Cat\,\mathrm{hot}$ is false by \auditref{C.P5}. Isomorphism of ind-objects alone does not supply either comparison; other path-model identifications retain \auditref{B.8}.\\
\addlinespace[3pt]
VIIb-C026 & Refuted generally; Proved under the stated criterion and for paths &
\auditref{A.4} refutes preorder cofinality plus parallel-map homotopies for arbitrary fibrations. \auditref{A.3} gives the exact homotopy-initiality criterion uniform over all diagrams. \auditref{C.P1}--\auditref{C.P3} prove the stronger path-specific descent argument and the actual comparison roof of \auditref{A.P2}.\\
\addlinespace[3pt]
VIIc-C004 & Translation Proved; faithfulness Refuted &
\auditref{E.1} proves that the literal even translation is the identity in $\Hot$. \auditref{E.13} independently verifies the finite-support counterexample to faithfulness of finite set-valued tests. The specified equality holds, while the proposed general inference from those tests is false.\\
\addlinespace[3pt]
VIIc-C005 & Proved for the eventual model; all-type strengthening Refuted &
For the eventual alternating endpoint functor, smoothness, properness, and the strict homotopy-fibre assertion are \auditref{E.1} and \auditref{E.4}. The all-type endpoint functor is smooth but need not be proper or a Grothendieck (op)fibration, by \auditref{E.3}. These are different functors.\\
\addlinespace[3pt]
\multicolumn{3}{@{}l@{}}{\textit{Numbered targets inherited from the earlier packages}}\\
\addlinespace[2pt]
CATHOT E.26(1) & Proved &
\auditref{C.P3}: for strictly commuting $u:t\to s$ between preorder-cofinal projective type sequences, the actual $u^*:C_\infty(s)\to C_\infty(t)$ is in $W_\infty$, for every small $M$, wide $\Sigma$, and endpoint pair.\\
\addlinespace[3pt]
CATHOT E.26(2) & Proved &
\auditref{C.P3}: the comparisons through $t*s$ equal $\beta_t^{-1}\beta_s$ in $\Hot$. This proves independence and transitivity for the specified maps and roofs, rather than merely an abstract isomorphism of their sources and targets.\\
\addlinespace[3pt]
CATHOT E.26(3) & Proved &
\auditref{C.P3}: the specified fibred Grothendieck functor $E_t\to\Ch_\Sigma(M;a,b)$ is in $W_\infty$. \auditref{A.P2} identifies its roof with the sequential colimit. The proof does not assume that $J$ is a weak equivalence.\\
\addlinespace[3pt]
CATHOT E.27 & Proved for every cofinal block &
\auditref{C.P4}: for any block containing both signs, both actual padding maps are weak equivalences, and the resulting operation is associative, unital, and independent of the block in $\Hot$. A one-sign block does not give a cofinal type system and is outside this assertion.\\
\addlinespace[3pt]
CATHOT E.17 & Proved in the small-index scope &
\auditref{D.K1}--\auditref{D.K4} prove the acyclic-map reduction, the cardinal-controlled Kan--Thurston input, and factorization through $\Pi_1$ for a small filtered ind-system tested on all small categories in the same universe. \auditref{D.K5} applies this to the countable path system and records the bounded-shape limitation.\\
\addlinespace[3pt]
STRATA E.11 & Open; standard path constructions Refuted as solutions &
\auditref{E.6} proves that the eventual endpoint construction fails both strict lifting properties over the five-element rectangular-band monoid; the all-type construction fails the opfibration property. \auditref{E.11} retains existence of a different strict bifibration realizing its path local system.\\
\addlinespace[3pt]
\multicolumn{3}{@{}l@{}}{\textit{The three leads and the additional rectification targets}}\\
\addlinespace[2pt]
Lead 1 & Known restricted model; Open general inclusion &
\auditref{B.1} specifies $J$ and \auditref{B.2} identifies its restricted source with the hammock mapping space. \auditref{B.3} confirms $J$ need not be a categorical isomorphism, even with saturated markings. The general weak-equivalence question, including the model and Brown cases, remains \auditref{B.8}.\\
\addlinespace[3pt]
Lead 2 & Proved, with a correction &
\auditref{B.4} proves the assertion whenever $\Sigma\subseteq\operatorname{Iso}(M)$, with a uniform contraction of the full category in $\Cat\,\mathrm{hot}$. The suggested initial object need not exist in the total category. For all arrows marked, $J$ is the identity.\\
\addlinespace[3pt]
Lead 3 & Proved sufficient implication; finite-example claim Refuted &
\auditref{D.7} proves the rectification statements with their actual finite-direct or universe-relative shape bounds; \auditref{D.4} proves the second (b). The claim that the manuscript's finite example fails rectification is false by \auditref{D.2}: it rectifies every shape and satisfies Kaledin's axioms.\\
\addlinespace[3pt]
Finite-shape recognition and product tests & Refuted in the cubical form; Open in the stated generality &
\auditref{D.10} refutes the proposed cube criterion. \auditref{D.21} upgrades the canonical comparison when all Kaledin product shapes are available. Necessary-and-sufficient criteria using only finite posets, finite categories, or a prescribed cardinal bound, and the finite-product-shape upgrade, remain \auditref{D.22}.\\
\addlinespace[3pt]
Intrinsic cofibrations and later axioms & Proved with the stated existence hypotheses; Open without them &
\auditref{D.8} proves finite rectification with an initial object and pushouts; \auditref{D.9} refutes automatic cofibrancy of every object. \auditref{D.17} compares XVII, XIX, and XIII and identifies the transfinite hypotheses for all-shape rectification. The version without an initial object remains \auditref{D.22}.\\
\addlinespace[3pt]
Barwick--Kan fibrancy; $B\mathbb Z$ & Two implications Refuted; finite-shape implication Open &
\auditref{D.11} gives a Barwick--Kan fibrant example rectifying finite posets but not $B\mathbb Z$. \auditref{D.12} gives an all-shape rectifiable example which is not fibrant. Whether fibrancy forces finite-poset or finite-category rectification remains \auditref{D.22}.\\
\addlinespace[3pt]
Relative poset presentations & Refuted as a rectification criterion &
\auditref{D.13} gives a relative poset presenting an ordinary category whose strict diagram family nevertheless fails the specified comparison. Presenting the correct localization does not by itself rectify diagram categories.\\
\addlinespace[3pt]
Kaledin recognition and extension & Known with actual shapes; naive extension obstruction Proved &
\auditref{D.14} audits the axioms and shape classes. \auditref{D.15} identifies the object represented by a naive enhanced family as its specified localization. \auditref{D.16} locates an explicit countable-poset failure of additivity; the finite example of \auditref{D.2} has no failed enhanced-category axiom.\\
\addlinespace[3pt]
Accessible localizers; descent classes & Known with accessibility; broader assertion Open &
\auditref{D.18} gives all-shape rectification and enough intrinsic cofibrations for accessible localizers. \auditref{D.23} proves that every nonzero spectrum's homology localizer lacks enough intrinsic fibrations. \auditref{D.19} refutes two-out-of-three for the stated descent class. The nonaccessible derivator question remains \auditref{D.22}.\\
\addlinespace[3pt]
Recovery from diagram homotopy categories & Known with the full coherent family &
\auditref{D.20} resolves the Arlin and Fuentes-Keuthan--K\k edziorek--Rovelli references and distinguishes their shape and morphism conventions. Its arrow-shape example refutes replacing coherent diagram homotopy categories by diagram categories in the homotopy category.\\
\end{longtable}
\endgroup
